\documentclass[11pt]{article}
\usepackage[a4paper,margin=25mm]{geometry}
\usepackage{amsmath,amssymb,amsthm,booktabs,array}
\usepackage[T1]{fontenc}
\usepackage{lmodern}
\usepackage[colorlinks=true,urlcolor=blue,citecolor=blue,linkcolor=blue]{hyperref}
\newcommand{\tr}{\operatorname{tr}}
\newcommand{\Id}{I}
\newcommand{\CH}{C_{\mathrm H}}
\newcommand{\CN}{C_{\mathrm{NH}}}
\newcommand{\CM}{C_{\mathrm{MI}}}
\newcommand{\normone}[1]{\left\lVert#1\right\rVert_1}
\newtheorem{theorem}{Theorem}
\newtheorem{lemma}[theorem]{Lemma}
\newtheorem{corollary}[theorem]{Corollary}
\title{Fixed-purity Holevo minima extending Das et al. (2025): a spectral proof of the full-ONB conjecture after Theorem 3 in Section 3.2.1}
\author{Haobo Ma\\\small ChronoAI Pte Ltd, Singapore\and
Wenlin Zhang\\\small National University of Singapore\\\small ChronoAI Pte Ltd, Singapore}
\date{November 8, 2026}
\begin{document}
\maketitle
\begin{abstract}
Extending Das et al., Quantum \textbf{9}, 1867 (2025), we prove their
full-orthonormal-basis tomography conjecture immediately after Theorem 3
in Section 3.2.1 by completing their reduction in Section 5.4 after Fig.~5:
the Nagaoka--Hayashi bound is at most $d+1$ times the Holevo bound once
the latter is minimized by the maximally mixed state.
For full-rank states the same sharp bound holds for the optimal
individual-measurement cost, improving the source's Theorem~5 as well as Theorem~3.
For identity weight and purity $P$, we prove the stronger estimate
$\CN/\CH\leq(d^2+d-1-P)/m_d(P)$, where
$m_d(P)=d-P+\sqrt{(d-1)(dP-1)}$ is the exact minimum Holevo cost at that purity.
We prove that minimality through an explicit spectral evaluation of the
Holevo cost, within the established D-invariant theory, and use the
feasible Nagaoka--Hayashi blocks of Das et al. The maximum is attained
at the maximally mixed state and is strictly smaller at every other state.
We also distinguish this increasing minimum cost from the failure of purity
to order the costs of individual states.
\end{abstract}

\section{Question and normalization}
Das et al.\ \cite[Sec.~3.2.1, paragraph after Theorem~3; Sec.~5.4,
paragraph after Eq.~(49)]{Das} propose replacing their $d+2$ bound by $d+1$
for full orthonormal-basis (ONB) tomography: the source says
``we propose'' $\mathcal R^{\mathrm{NH}}[\{\rho_\theta\}]\leq d+1$.
In Section~5.4, immediately after Fig.~5, they also identify the sufficient
condition $\CH(\rho)\geq\CH(\Id_d/d)$ for this improvement.
Our contribution is a proof of that minimality by spectral evaluation,
its exact fixed-purity refinement, and an individual-measurement extension.
We use the identity weight. The first arXiv version (15 May 2024), Section~IV.B
after Theorem~2 and Eq.~(32), already identifies the same $d+1$ target;
the locators used below refer to version~2 (25 September 2025).
For full-rank full tomography, the equality of the Holevo and right
logarithmic derivative bounds belongs to the established D-invariant
theory; see Suzuki~\cite[Sec.~4.1.2, Proposition~2]{Suzuki}.
The eigenvalue-gap expression below is a direct evaluation in the present
normalization. Its proof also includes positive semidefinite boundary states.
Scott's linear-tomography error formula~\cite[Sec.~VI, Eqs.~(101)--(103)]{Scott}
is exactly the score-variance formula used below: Hilbert--Schmidt squared
error equals the sum of squared errors in these ONB coordinates.
The individual-measurement corollary combines that prior estimator bound
with the spectral Holevo minimum.
Fix $d\geq2$, $n=d^2-1$, and a basis
$B_1,\ldots,B_n$ of the traceless Hermitian $d\times d$ matrices with
$\tr(B_jB_k)=\delta_{jk}$. Write
\[
 \rho=\frac{\Id_d}{d}+\sum_{j=1}^n\theta_j B_j\succeq0,
 \qquad P=\tr(\rho^2)=\frac1d+\sum_j\theta_j^2.
\]
Thus $\tr\rho=1$ and $1/d\leq P\leq1$. These are the normalization and
model of \cite[Sec.~5.2, Eq.~(23); Appendix~A]{Das}.

An uncentered locally unbiased family consists of Hermitian operators
$X_k$ satisfying
\[
 \tr(\rho X_k)=\theta_k,\qquad
 \tr(B_jX_k)=\delta_{jk}.
\]
Every such family equals $X_k=B_k$: the second constraint gives
$X_k=B_k+a_k\Id_d$, and the first gives $a_k=0$.
For $Z_{jk}=\tr(\rho B_jB_k)$, the identity-weight Holevo bound is therefore
\begin{equation}\label{eq:costs}
 \CH=\tr Z+\normone{\operatorname{Im}Z}-\sum_j\theta_j^2.
\end{equation}
Here $\operatorname{Im}$ is entrywise and $\normone{\cdot}$ is the sum of
singular values. The Nagaoka--Hayashi bound is
\begin{equation}\label{eq:nh}
 \CN=\inf_L\sum_j\tr(\rho L_{jj})-\sum_j\theta_j^2,
 \quad L_{jk}=L_{kj}=L_{jk}^{\dagger},\quad
 [L_{jk}-B_jB_k]_{j,k}\succeq0.
\end{equation}
The block matrix acts on $\mathbb C^n\otimes\mathbb C^d$.
These are \cite[Eqs.~(4), (6)]{Das}; the semidefinite formulation is
developed in \cite{Conlon}. Singular $\rho$ are included as algebraic
values of these formulas; the usual full-rank statistical interpretation
applies on the interior.

\section{The estimate}
\begin{theorem}\label{thm:main}
For every $d\geq2$, every full traceless Hermitian ONB, and every density
matrix $\rho$ as above,
\begin{equation}\label{eq:main}
 \CH\geq(d-1)(1+P),\qquad
 \CN\leq d^2+d-1-P,
\end{equation}
and consequently
\begin{equation}\label{eq:ratio}
 \frac{\CN}{\CH}
 \leq \frac{d^2+d-1-P}{(d-1)(1+P)}\leq d+1.
\end{equation}
The second inequality is strict unless $\rho=\Id_d/d$.
At $\rho=\Id_d/d$ the ratio is $d+1$.
\end{theorem}

\begin{lemma}\label{lem:spectrum}
If $\lambda_1,\ldots,\lambda_d$ are the eigenvalues of $\rho$, then
\begin{equation}\label{eq:spectrum}
 \CH=d-P+\sum_{a<b}|\lambda_a-\lambda_b|.
\end{equation}
In particular $\CH\geq(d-1)(1+P)$.
\end{lemma}
\begin{proof}
Adjoining $\Id_d/\sqrt d$ gives an orthonormal basis for all complex
matrices. Expansion in that basis, evaluated on matrix units, gives
\begin{equation}\label{eq:fierz}
 \sum_j(B_j)_{ab}(B_j)_{ce}
 =\delta_{ae}\delta_{bc}-\frac1d\delta_{ab}\delta_{ce},\qquad
 \sum_j B_j A B_j=\tr(A)\Id_d-\frac{A}{d}.
\end{equation}
These are the usual completeness identities; see
\cite[Sec.~2]{Haber}, with generators multiplied by $\sqrt2$.
Putting $A=\Id_d$ yields $\sum_jB_j^2=(d^2-1)\Id_d/d$.
Hence $\tr Z=(d^2-1)/d$, and subtracting
$\sum_j\theta_j^2=P-1/d$ leaves $d-P$ in \eqref{eq:costs}.

An orthogonal change of the real Hermitian ONB changes $Z$ to $OZO^T$,
which preserves its imaginary-part trace norm. Conjugating both $\rho$
and the basis by a unitary also preserves every trace. We may therefore
diagonalize $\rho$ and use, for each $a<b$,
\[
 S_{ab}=\frac{E_{ab}+E_{ba}}{\sqrt2},\qquad
 A_{ab}=\frac{-i(E_{ab}-E_{ba})}{\sqrt2},
\]
together with an orthonormal traceless diagonal family.
The imaginary part of $Z$ has one block per pair $a<b$,
\[
 \begin{pmatrix}0&(\lambda_a-\lambda_b)/2\\
 -(\lambda_a-\lambda_b)/2&0\end{pmatrix},
\]
and zeros elsewhere. Indeed two distinct unordered off-diagonal pairs
have no diagonal contribution to their product, diagonal/off-diagonal
products have zero trace against $\rho$, and diagonal/diagonal products
are real. The two singular values of the displayed block are both
$|\lambda_a-\lambda_b|/2$. This proves \eqref{eq:spectrum}.

Since $0\leq\lambda_a\leq1$,
\[
 \sum_{a<b}|\lambda_a-\lambda_b|
 \geq\sum_{a<b}(\lambda_a-\lambda_b)^2=dP-1,
\]
where the last equality follows by expanding and using
$\sum_a\lambda_a=1$. Substitution in \eqref{eq:spectrum} proves the claim.
\end{proof}

\begin{proof}[Proof of Theorem~\ref{thm:main}]
The first bound is Lemma~\ref{lem:spectrum}.
For the second use the blocks of Das et al.\ \cite[Eq.~(37), Appendix~E.2]{Das}:
\begin{equation}\label{eq:primal}
 L^*_{jk}=c(B_jB_k+B_kB_j+\delta_{jk}\Id_d),
 \qquad c=\frac{d+1}{d+2}.
\end{equation}
Here is a direct positivity argument. Put $r=1/(d+1)$ and define the
Hermitian block matrix $T_{jk}=B_kB_j-rB_jB_k$. By \eqref{eq:fierz},
\begin{align*}
 (T^2)_{j\ell}
 &=\delta_{j\ell}\Id_d+
 \left(-\frac1d+\frac{2r}{d}+\frac{r^2(d^2-1)}d\right)B_jB_\ell
 =\delta_{j\ell}\Id_d.
\end{align*}
Thus the eigenvalues of $T$ belong to $\{-1,1\}$ and $\Id_{nd}+T\succeq0$.
Since $cr=1-c$, the constraint matrix for \eqref{eq:primal} is
$[L^*_{jk}-B_jB_k]=c(\Id_{nd}+T)\succeq0$.
The other constraints in \eqref{eq:nh} follow directly from
$B_j=B_j^{\dagger}$. Also,
\[
 \sum_j L^*_{jj}
 =c\left(2\frac{d^2-1}{d}+d^2-1\right)\Id_d
 =\frac{(d^2-1)(d+1)}d\Id_d.
\]
The infimum is finite: every admissible $L$ has
$L_{jj}\succeq B_j^2\succeq0$, so its uncentered cost is nonnegative;
$L^*$ supplies a finite feasible cost. Inserting $L^*$ in
\eqref{eq:nh} and subtracting $P-1/d$ proves the second bound.

The denominator $(d-1)(1+P)$ is positive. The difference between
$(d+1)(d-1)(1+P)$ and $d^2+d-1-P$ is $d(dP-1)$, proving
\eqref{eq:ratio}, strictly for $P>1/d$. At $P=1/d$,
$\rho=\Id_d/d$ and $\CH=(d^2-1)/d$ by \eqref{eq:spectrum}.
The matching value $\CN=(d^2-1)(d+1)/d$ is
\cite[Lemma~3, Eq.~(38), Appendix~E.3--E.4]{Das}, whose matching primal
and dual objectives establish equality without any SIC-existence assumption.
\end{proof}

\section{Consequences for purity and the source's related statements}
\begin{corollary}\label{cor:depolarized}
For $\rho=p|\psi\rangle\langle\psi|+(1-p)\Id_d/d$, $0\leq p\leq1$,
\[
 \CH=\frac{d^2-1}{d}+(d-1)p-\frac{d-1}{d}p^2.
\]
For pure states, $\CH=2(d-1)$. The maximally mixed state is the unique
minimizer of $\CH$ over all density matrices in this model.
\end{corollary}
\begin{proof}
The eigenvalues are $(1+(d-1)p)/d$ and $(1-p)/d$ with multiplicity $d-1$.
Their absolute pairwise differences sum to $(d-1)p$, and
$P=(1+(d-1)p^2)/d$. Apply \eqref{eq:spectrum}.
The final assertion follows from
$\CH\geq(d-1)(1+P)\geq(d^2-1)/d$ and equality only at $P=1/d$.
\end{proof}

\begin{corollary}\label{cor:fixedpurity}
At each $P\in[1/d,1]$, the minimum Holevo cost over states of purity $P$ is
\begin{equation}\label{eq:envelope}
 m_d(P)=d-P+\sqrt{(d-1)(dP-1)}.
\end{equation}
It is attained by the depolarized pure states of
Corollary~\ref{cor:depolarized}, with $p=\sqrt{(dP-1)/(d-1)}$.
For $P>1/d$, the other minimizing spectra, when distinct, have one smaller
eigenvalue and $d-1$ equal larger eigenvalues; they exist exactly when
$P\leq1/(d-1)$. Consequently,
\begin{equation}\label{eq:sharpratio}
 \frac{\CN}{\CH}\leq F_d(P):=
 \frac{d^2+d-1-P}{m_d(P)}
 \leq \frac{d^2+d-1-P}{(d-1)(1+P)}\leq d+1.
\end{equation}
The function $F_d$ is strictly decreasing on $[1/d,1]$.
\end{corollary}
\begin{proof}
Order the eigenvalues $\lambda_1\geq\cdots\geq\lambda_d$ and put
$\delta_k=\lambda_k-\lambda_{k+1}\geq0$.
Let $v^{(k)}$ have $k$ coordinates equal to $(d-k)/d$ and the remaining
coordinates equal to $-k/d$. Then
\[
 \lambda-(1/d,\ldots,1/d)=\sum_{k=1}^{d-1}\delta_kv^{(k)},\quad
 \|v^{(k)}\|_2=\sqrt{\frac{k(d-k)}d},\quad
 \sum_{a<b}|\lambda_a-\lambda_b|=\sum_{k=1}^{d-1}k(d-k)\delta_k.
\]
The triangle inequality and $k(d-k)\geq d-1$ give
\[
 \sqrt{P-1/d}\leq\sum_k\delta_k\sqrt{k(d-k)/d}
 \leq\frac{\sum_{a<b}|\lambda_a-\lambda_b|}{\sqrt{d(d-1)}}.
\]
Equation~\eqref{eq:spectrum} now proves $\CH\geq m_d(P)$, and the
depolarized spectrum gives equality. Equality in the last two inequalities,
away from the uniform spectrum, requires exactly one nonzero gap, at
$k=1$ or $k=d-1$: different $v^{(k)}$ are not positively collinear.
The second spectrum is
$( (1+q)/d,\ldots,(1+q)/d,(1-(d-1)q)/d )$ with
$q=\sqrt{(dP-1)/(d-1)}\leq1/(d-1)$, proving the equality description.
The two families coincide up to permutation when $d=2$.
Since $0\leq dP-1\leq d-1$, we have
$\sqrt{(d-1)(dP-1)}\geq dP-1$, so
$m_d(P)\geq(d-1)(1+P)$.
For $P>1/d$,
\[
 m_d'(P)=-1+\frac{d(d-1)}{2\sqrt{(d-1)(dP-1)}}\geq d/2-1\geq0.
\]
The derivative is positive except possibly at the endpoint $d=2,P=1$.
Thus $m_d$ is strictly increasing; the positive numerator of $F_d$ is
strictly decreasing. Continuity includes $P=1/d$ and proves the last claim.
\end{proof}

\begin{corollary}\label{cor:mi}
For full identity-weight ONB tomography of a full-rank state, let $\CM$
denote the optimal individual-measurement cost, as in \cite[Eq.~(3)]{Das}.
Then
\[
 \CN\leq\CM\leq d^2+d-1-P,\qquad
 \frac{\CM}{\CH}\leq F_d(P)\leq d+1.
\]
The maximum ratio is $d+1$, attained exactly at $\rho=\Id_d/d$.
No existence assumption for SIC POVMs is needed.
The Gill--Massar ratio $C_{\mathrm{GM}}/\CH$ has the same upper envelope
$F_d(P)$ and the same unique maximizer.
\end{corollary}
\begin{proof}
The upper bound has the same right side as \cite[Eq.~(55)]{Das}.
Here is a finite-POVM argument, using weighted complex projective
two-designs and the linear-estimation variance formula of
\cite[Proposition~13, Eqs.~(101)--(103)]{Scott}.
For a uniformly distributed rank-one projector $Q$,
unitary invariance and trace normalization give
\[
 \int Q\otimes Q\,d\mu(Q)=\frac{\Id+\mathsf F}{d(d+1)},
\]
where $\mathsf F$ exchanges the tensor factors. The integral belongs to
the convex hull of the compact set $\{Q\otimes Q\}$, so Carath\'eodory's
theorem gives finitely many projectors $Q_s$ and weights $w_s>0$ with
$\sum_s w_s=1$ and
\begin{equation}\label{eq:design}
 \sum_s w_s Q_s\otimes Q_s=\frac{\Id+\mathsf F}{d(d+1)}.
\end{equation}
Taking a partial trace gives $\sum_s w_sQ_s=\Id_d/d$.
Thus $M_s=dw_sQ_s$ is a POVM. Assign the outcome scores
$f_j(s)=(d+1)\tr(B_jQ_s)$. Equation~\eqref{eq:design} gives
\[
 \sum_s\tr(\rho M_s) f_j(s)=\tr(\rho B_j)=\theta_j.
\]
Hence the scores are globally unbiased. Completeness of the ONB gives,
for every rank-one $Q_s$,
\[
 \sum_j f_j(s)^2=(d+1)^2\left(\tr(Q_s^2)-\frac1d\right)
 =\frac{(d^2-1)(d+1)}d.
\]
Their total variance is therefore
$(d^2-1)(d+1)/d-(P-1/d)=d^2+d-1-P$.
The moment identity also implies informational completeness: if traceless
Hermitian $A$ has $\tr(AQ_s)=0$ for every $s$, then
$\tr(A^2)/[d(d+1)]=0$ by \eqref{eq:design}, so $A=0$.
For full-rank $\rho$, the outcome probabilities are positive and the
classical Fisher matrix $J_M$ is nonsingular. The classical Cram\'er--Rao
inequality gives $\tr J_M^{-1}\leq\sum_j\operatorname{Var}f_j$, proving
the cost bound. The variance formula itself holds for every positive
semidefinite state.
The inequality $\CN\leq\CM$ is the NH relaxation
\cite[Eq.~(8)]{Das}; Corollary~\ref{cor:fixedpurity} proves the ratio bound.
At $\Id_d/d$, the matching NH value from Lemma~3 of \cite{Das} and the
displayed upper bound give equality for $\CM$ too. Away from that state,
$F_d(P)<d+1$.
The GM bound is a lower bound on $\CM$ \cite[Eq.~(158)]{Das}, so it
inherits the envelope. At $\Id_d/d$, the SLD Fisher matrix is
$d\Id_n$, giving $C_{\mathrm{GM}}=n^2/[d(d-1)]=(d^2-1)(d+1)/d$,
which proves its sharpness as well.
\end{proof}

\begin{corollary}\label{cor:range}
The Holevo cost satisfies
\[
 \frac{d^2-1}{d}\leq\CH\leq2(d-1).
\]
The lower endpoint occurs exactly at $\rho=\Id_d/d$, and the upper endpoint
exactly at pure states. Nevertheless, greater purity need not give greater
Holevo cost, even for full-rank states in dimension three.
\end{corollary}
\begin{proof}
The lower endpoint was proved in Corollary~\ref{cor:depolarized}.
Write $M=\sum_{a<b}\min(\lambda_a,\lambda_b)$.
Since $\sum_{a<b}|\lambda_a-\lambda_b|=d-1-2M$, we have
$\CH=2d-1-P-2M$.
Also $1-P=2\sum_{a<b}\lambda_a\lambda_b\leq2M$, giving the upper bound.
If two eigenvalues are positive, their product is strictly smaller than
their minimum, so equality is possible exactly at rank one.
For the final assertion take the two diagonal spectra
\[
 \lambda^{A}=(3/4,1/5,1/20),\qquad
 \lambda^{B}=(4/5,1/10,1/10).
\]
Both have gap sum $7/5$, while
$P_A=121/200<33/50=P_B$ and
$\CH(\lambda^A)=759/200>187/50=\CH(\lambda^B)$.
\end{proof}

The source's increasing-Holevo-cost discussion following Eq.~(49) in
\cite[Sec.~5.4]{Das} is supported there by particular families and numerical
plots. A universal ordering of arbitrary states by purity is not asserted
as a theorem in the source and fails for the displayed pair.
The exact increasing minimum in Corollary~\ref{cor:fixedpurity}
is a valid state-independent statement. Its decreasing ratio bound $F_d(P)$ does not by
itself determine the actual maximal NH/H ratio at fixed purity, nor its
maximizers in \cite[Sec.~3.3, Fig.~3]{Das}.

\begin{corollary}\label{cor:stability}
Let $R$ be either $\CN/\CH$, or $\CM/\CH$ for full-rank states. Then
\[
 d+1-R\geq\frac{d(dP-1)}{(d-1)(1+P)}.
\]
In particular, $R\geq d+1-\varepsilon$ implies
$\|\rho-\Id_d/d\|_{\mathrm{HS}}^2\leq2(d-1)\varepsilon/d^2$.
\end{corollary}
\begin{proof}
For either ratio, $R\leq(d^2+d-1-P)/[(d-1)(1+P)]$.
Subtract this bound from $d+1$.
Then use $P-1/d=\|\rho-\Id_d/d\|_{\mathrm{HS}}^2$ and $1+P\leq2$.
\end{proof}

\section{Coverage and mathematical evidence}
The following locators refer to \cite{Das}, arXiv version~2.
\begin{center}\small
\begin{tabular}{@{}p{.43\linewidth}p{.51\linewidth}@{}}
\toprule
Source item & Coverage\\\midrule
Theorem~1; Eq.~(17); Appendix~J: general $n$ and GM ratio bounds
& Prior proofs cited. For full ONB tomography, Corollary~\ref{cor:mi} also gives the GM purity envelope; NH--GM equality is not needed.\\[3pt]
Sec.~3.2.1 after Theorem~3; Sec.~5.4 after Eq.~(49): full-ONB $d+1$ conjecture
& Proved: Theorem~\ref{thm:main}, completing the reduction already identified by Das et al. (2025), Sec.~5.4 after Fig.~5.\\[3pt]
Theorem~3: full-ONB $d+2$ estimate
& Proved: improved to $d+1$; equality only at $\Id_d/d$.\\[3pt]
Sec.~5.4 after Fig.~5: minimum-Holevo-cost reduction
& Prior reduction credited to Das et al. (2025); minimality proved by Lemma~\ref{lem:spectrum}. Exact fixed-purity minimum and global range in Corollaries~\ref{cor:fixedpurity}--\ref{cor:range}.\\[3pt]
Sec.~5.4 after Eq.~(49): increasing-Holevo-cost language
& Universal ordering not asserted as a theorem in the source; it fails for the pair in Corollary~\ref{cor:range}. The fixed-purity minimum increases.\\[3pt]
Theorem~2, Lemma~3: maximally mixed ratio
& Prior proof: Das et al. (2025), Lemma~3 and Appendix~E.3--E.4, supplies sharpness.\\[3pt]
Eq.~(18): depolarized pure-state Holevo cost
& Reproved; prior formula: Das et al. (2025), Eq.~(18), in Corollary~\ref{cor:depolarized}.\\[3pt]
Eq.~(19), fixed-purity extremizers in Sec.~3.3
& Not covered: NH-cost formula and actual ratio extremizers. Corollary~\ref{cor:fixedpurity} supplies a decreasing upper envelope, not their solution.\\[3pt]
Eqs.~(13), (50); Appendix~H
& Not covered: arbitrary-weight, arbitrary-state conjecture. Prior proof cited for the maximally mixed case.\\[3pt]
Theorem~4; general $n$-parameter conjecture in Sec.~4
& Prior proof cited for the maximally mixed subset bound; not covered for arbitrary subsets or the general $\min(n,d+1)$ conjecture.\\[3pt]
Theorem~5; Sec.~5.6, Eq.~(55); Appendix~I: MICRB $d+2$ bound
& Proved: improved to $d+1$ in Corollary~\ref{cor:mi}, with the same purity envelope and unique maximizer.\\[3pt]
Scott (2006), Sec.~VI, Eqs.~(101)--(103): two-design estimator variance
& Reproved; prior proof Scott (2006). The ONB-coordinate variance is the same Hilbert--Schmidt error; Corollary~\ref{cor:mi} combines it with the Holevo minimum.\\[3pt]
Sec.~4: qutrit four-parameter optimum
& Not covered: the optimum for general four-parameter models.\\
\bottomrule
\end{tabular}
\end{center}

The declaration
\texttt{D5.S3.Quantum.Measurements.}\allowbreak\texttt{OnbTomographyHolevoRatio.result}
in the \href{https://github.com/the-omega-institute/trureturing/blob/3d88df0f99ef3c8bfb034fa52be669b13746a844/D5/S3/Quantum/Measurements/OnbTomographyHolevoRatio.lean}{pinned Lean file}
proves $\CN\leq(d+1)\CH$ under the full identity-weight hypotheses of
Theorem~\ref{thm:main}, including positive semidefinite states.
The spectral formula, stronger ratio bounds, exact fixed-purity minimum
and equality spectra, MI bound, global range, counterexample and stability
have written proofs above. Sharpness uses the cited maximally mixed NH value.
The exact finite checks in \texttt{verify.py} cover matrix identities for
$d=2,3$; all $2379$ denominator-$12$ spectra for $2\leq d\leq5$;
$182$ attaining spectra for $2\leq d\leq8$; and two-design moment,
unbiasedness and Fisher identities for $6$- and $12$-outcome POVMs.
They complement the all-dimensional proofs. The stronger ratio envelope
need not be attained away from $\Id_d/d$.

\begingroup\small
\begin{thebibliography}{9}
\bibitem{Das}
A. Das, L. O. Conlon, J. Suzuki, S. K. Yung, P. K. Lam, and S. M. Assad,
\emph{Holevo Cram\'er--Rao bound: How close can we get without entangling measurements?},
Quantum \textbf{9}, 1867 (2025).
\href{https://doi.org/10.22331/q-2025-09-30-1867}{doi:10.22331/q-2025-09-30-1867};
\href{https://arxiv.org/abs/2405.09622v2}{arXiv:2405.09622v2}.
The earlier statement is in
\href{https://arxiv.org/abs/2405.09622v1}{arXiv:2405.09622v1}, Sec.~IV.B.
\bibitem{Conlon}
L. O. Conlon, J. Suzuki, P. K. Lam, and S. M. Assad,
\emph{Efficient computation of the Nagaoka--Hayashi bound for multiparameter estimation with separable measurements},
npj Quantum Information \textbf{7}, 110 (2021).
\href{https://doi.org/10.1038/s41534-021-00414-1}{doi:10.1038/s41534-021-00414-1}.
\bibitem{Haber}
H. E. Haber,
\emph{Useful relations among the generators in the defining and adjoint representations of SU($N$)},
SciPost Physics Lecture Notes \textbf{21} (2021).
\href{https://doi.org/10.21468/SciPostPhysLectNotes.21}{doi:10.21468/SciPostPhysLectNotes.21};
\href{https://arxiv.org/abs/1912.13302}{arXiv:1912.13302}.
\bibitem{Scott}
A. J. Scott,
\emph{Tight informationally complete quantum measurements},
Journal of Physics A: Mathematical and General \textbf{39}, 13507--13530 (2006).
\href{https://doi.org/10.1088/0305-4470/39/43/009}{doi:10.1088/0305-4470/39/43/009};
\href{https://arxiv.org/abs/quant-ph/0604049}{arXiv:quant-ph/0604049}.
\bibitem{Suzuki}
J. Suzuki,
\emph{Information Geometrical Characterization of Quantum Statistical Models in Quantum Estimation Theory},
Entropy \textbf{21}, 703 (2019).
\href{https://doi.org/10.3390/e21070703}{doi:10.3390/e21070703};
\href{https://arxiv.org/abs/1807.06990}{arXiv:1807.06990}.
The preprint's corresponding locators are Theorems~III.6 and IV.2 and Eq.~(31).
\end{thebibliography}
\endgroup
\smallskip\noindent\small
Licensed under Creative Commons Attribution 4.0 International (CC BY 4.0).
\end{document}
